\section{The relative Iitaka reduction and its Hodge line}
\label{setup:section}

The first task is to replace the fibers by their Iitaka images while
retaining every sufficiently divisible pluricanonical system. We first
fix the birational models on which divisorial order tests suffice. The
root form of a Kodaira-zero fiber will then provide a Hodge line and an
exact comparison of sections on the intermediate base.

We use additive notation for rational line bundles; equality of them
means equality in $\operatorname{Pic}\otimes\Q$. A rational line is
effective if a positive integral multiple has a nonzero section.
All degree assertions below are made after clearing denominators.
For a fibration $g:(Z,\Gamma)\to T$ and a very general smooth
fiber $F$, adjunction gives
\begin{equation}\label{setup:fiber-adjunction}
 (K_Z+\Gamma)|_F=K_F+\Gamma|_F.
\end{equation}

\subsection{The fixed smooth models}

A smooth-base fibration $g:(Z,\Gamma)\to T$ is \emph{neat} if there is a
proper bimeromorphic orbifold morphism
$w:(Z,\Gamma)\to(Z_0,\Gamma_0)$ to a smooth pair, with
$w_*\Gamma=\Gamma_0$, such that every prime divisor sent by $g$ into
codimension at least two is also sent by $w$ into codimension at least
two. Campana's construction gives suitable neat models, model
independence, and the formula
\begin{equation}\label{setup:neat-formula}
 \kappa(g\mid\Gamma)=\kappa\bigl(T,K_T+B(g,\Gamma)\bigr)
 \quad\text{on a neat model}
\end{equation}
\cite[Definition~3.8, Proposition~3.10, Corollary~4.11, and
Remark~6.2]{CampanaOrbifolds}.
For the proof we only need the defining inequality, valid on every
allowed smooth model:
\begin{equation}\label{setup:base-comparison}
 \kappa\bigl(T,K_T+B(g,\Gamma)\bigr)
 \ \geq\ \kappa(g\mid\Gamma).
\end{equation}

\subsection{Reduced exceptional boundaries}

\begin{lemma}[Logarithmic modifications]\label{setup:log-modification}
Let $p:Z'\to Z$ be a proper bimeromorphic holomorphic map between smooth
compact manifolds. Suppose that $\Gamma$ has simple normal crossing
support and coefficients in $\Q\cap[0,1]$, and set
\[
 \Gamma'=p_*^{-1}\Gamma+\operatorname{Exc}(p)_{\mathrm{red}}.
\]
Assume, after further resolution if necessary, that this boundary has
simple normal crossing support. There is an effective $p$-exceptional
rational divisor $A$ such that
\begin{equation}\label{setup:log-discrepancy}
 K_{Z'}+\Gamma'=p^*(K_Z+\Gamma)+A.
\end{equation}
For every sufficiently divisible positive integer $m$, the natural
identification of meromorphic pluricanonical forms induces
\begin{equation}\label{setup:section-invariance}
 p_*\mathcal O_{Z'}\bigl(m(K_{Z'}+\Gamma')\bigr)
 =\mathcal O_Z\bigl(m(K_Z+\Gamma)\bigr).
\end{equation}
These identifications respect multiplication, so the divisible section
rings and their Iitaka dimensions agree. Moreover, $p$ satisfies
\eqref{setup:orbifold-morphism}.
\end{lemma}

\begin{proof}
The difference in \eqref{setup:log-discrepancy} is exceptional, since
$p$ is an isomorphism at the generic point of every target prime.
To compute its coefficient at an exceptional prime $E$, choose
compatible local canonical forms.
Near a general point of its center on $Z$, let $z_1,\ldots,z_n$ be
coordinates in which the boundary components containing that center are
$z_i=0$ for $1\leq i\leq r$, with coefficients $\gamma_i$.
Put $t_i=\ord_E(p^*z_i)$ and let $k_E$ be the order of the Jacobian of
$p$ along $E$. The pullback of
\[
 \frac{dz_1}{z_1}\wedge\cdots\wedge\frac{dz_r}{z_r}
 \wedge dz_{r+1}\wedge\cdots\wedge dz_n
\]
has at most a simple pole along $E$. Indeed, at a general smooth point of
$E$, each logarithmic factor has the form $t_i\,du/u$ plus a holomorphic
one-form, and a nonzero wedge product contains at most one factor
$du/u$. Consequently $k_E+1\geq\sum_i t_i$. The coefficient in question
is therefore
\[
 k_E+1-\sum_i\gamma_i t_i
 \ \geq\ \sum_i(1-\gamma_i)t_i\ \geq\ 0.
\]
For an effective integral exceptional divisor $mA$, normality of $Z$
gives $p_*\mathcal O_{Z'}(mA)=\mathcal O_Z$: a meromorphic function
with poles only on $mA$ descends outside the codimension-two exceptional
image and extends by Hartogs' theorem. The projection formula proves
\eqref{setup:section-invariance}, compatibly with products and ratios.

Finally, an exceptional prime has infinite boundary multiplicity.
A nonexceptional prime is a strict transform, with pullback order one
and unchanged coefficient. These two cases prove
\eqref{setup:orbifold-morphism}.
\end{proof}

\begin{corollary}[Descent across exceptional centers]
\label{setup:exceptional-descent}
In the situation of \Cref{setup:log-modification}, a meromorphic
$m$-pluricanonical form on $Z'$ satisfying the boundary bounds of
$m\Gamma'$ at every nonexceptional prime descends to a section of
$m(K_Z+\Gamma)$ whenever $m\Gamma$ and $m\Gamma'$ are integral.
The assertion also holds after tensoring by the pullback of a line
bundle on $Z$.
\end{corollary}

\begin{proof}
On the complement of the exceptional image the form is a holomorphic
section of the prescribed line bundle on $Z$: divisorial regularity
suffices on a smooth space. A local trivialization reduces extension
across the remaining codimension-two analytic subset to Hartogs'
extension theorem. The tensor-twisted assertion has the same proof.
\end{proof}

\begin{corollary}[Fibers, composition, and a fixed base]
\label{setup:modification-compatibility}
Let $p:(Z',\Gamma')\to(Z,\Gamma)$ be a modification with the boundary
rule of \Cref{setup:log-modification}. Then:
\begin{enumerate}[label=\textup{(\roman*)}]
\item If $p$ is a modification over a fixed base, it induces a
logarithmic modification on a very general smooth fiber. The log
Kodaira dimensions of the corresponding fibers agree.
\item A finite composition of such modifications again has boundary
equal to the strict transform of the initial boundary plus the reduced
exceptional divisor of the composite.
\item If $g:(Z,\Gamma)\to T$ has smooth base and $g'=g\circ p$, then
\begin{equation}\label{setup:fixed-base-boundary}
 B(g',\Gamma')=B(g,\Gamma).
\end{equation}
\end{enumerate}
\end{corollary}

\begin{proof}
For (i), choose the parameter so that the relevant smoothness and
dimension statements hold for the finitely many exceptional components
and their images. A component dominating the base restricts to a
reduced divisor, and its exceptional image restricts to codimension at
least two. Components not dominating the base do not meet this fiber.
The restricted map has the same boundary rule, so
\Cref{setup:log-modification} applies on the fiber. For (ii), every
prime exceptional for the composite is either newly exceptional or the
transform of an earlier exceptional prime; in both cases its coefficient
is one.

For (iii), fix a prime $D\subset T$. Every old prime dominating $D$
has a strict transform, with the same pullback order and boundary
multiplicity. Every new prime in the minimum defining
$m(g',\Gamma';D)$ is exceptional for $p$, so its boundary multiplicity
is infinite. Adding these entries leaves the old minimum unchanged,
including when that minimum was already infinite.
\end{proof}

For choice independence in the singular-base definition, take a common
resolution of two chosen diagrams. Their induced source boundaries
coincide: old primes retain their coefficients and all primes exceptional
over the initial $X$ have coefficient one. The comparison maps satisfy
\Cref{setup:log-modification}, so the smooth fibrations are equivalent
and their total-space and very general fiber log Kodaira dimensions agree.

\subsection{Flattening relative to a fixed reference pair}

To compare forms over generic points of base divisors, we need control of
source divisors over higher codimension. Proper analytic flattening gives
a base modification whose strict transform is flat
\cite{HironakaFlattening,FujikiDouady}. Here the strict transform is the
closure of the original family over the isomorphism locus, with base
torsion removed. Further base change preserves flatness. We use only
the resulting equidimensionality before resolving this intermediate
source.

\begin{lemma}[Flattening and the exceptional reference]
\label{setup:flattening}
Let $g_0:(Z_0,\Gamma_0)\to T_0$ be a fibration from a smooth compact
pair with rational SNC boundary in $[0,1]$ to a smooth compact manifold.
There are a smooth-base modification
$q:T\to T_0$, a smooth-source modification $p:Z\to Z_0$, and a
fibration $g:Z\to T$ with $q\circ g=g_0\circ p$ such that, for
\[
 \Gamma=p_*^{-1}\Gamma_0+\operatorname{Exc}(p)_{\mathrm{red}},
\]
the boundary has simple normal crossing support and every prime divisor
$E\subset Z$ satisfying $\operatorname{codim}_Tg(E)\geq2$ is
$p$-exceptional. In particular every such prime has coefficient one in
$\Gamma$, and the resulting fibration is neat with reference pair
$(Z_0,\Gamma_0)$.

One may resolve prescribed proper analytic bad loci on the base during
this construction. After any finite sequence of further base
modifications, one can choose dominating smooth source models for which
the same exceptional assertion holds with respect to the original
reference $Z_0$.
\end{lemma}

\begin{proof}
Flatten first, before resolving the source. After resolving the new base
by further base change, denote the equidimensional strict transform by
$\bar g:\bar Z\to T$, and its modification to $Z_0$ by $\bar p$.
Resolve $\bar Z$ and all boundary data, and write
\[
 Z\xrightarrow{r}\bar Z\xrightarrow{\bar p}Z_0,
 \qquad p=\bar p\circ r,
 \qquad g=\bar g\circ r.
\]
Let $N=\dim Z_0$, $b=\dim T$, and $d=N-b$. If a prime $E$ on $Z$
has image contained in an analytic subset $C\subset T$ of codimension
at least two, equidimensionality gives
\[
 \dim\bar g^{-1}(C)\leq\dim C+d\leq N-2.
\]
It follows that $r(E)$, and therefore also $p(E)$, has dimension at most
$N-2$. Thus $E$ is exceptional over $Z_0$. The boundary and
orbifold-morphism assertions follow from \Cref{setup:log-modification}.
The new map has connected general fiber.
Its Stein factorization has a finite bimeromorphic map to the normal
base $T$, which is an isomorphism, so all its fibers are connected.

For any further base modification $T'\to T$, pull back the flat
intermediate family and take a common resolution dominating the previous
smooth source. The same dimension estimate proves exceptionality over
$Z_0$. This includes principalizations of prescribed bad loci and their
SNC resolutions on the base. Iterating proves the persistence assertion;
\Cref{setup:modification-compatibility}\textup{(ii)} identifies the
stepwise boundary with that computed directly over $Z_0$.
\end{proof}

The resolution need not be flat: the conclusion we retain is
exceptionality over $Z_0$, which permits
\Cref{setup:exceptional-descent}.

All constructions remain available in Fujiki class $\mathcal C$.
A smooth compact member admits a compact K\"ahler modification
\cite[Theorems~0.7 and~3.4]{DemaillyPaun2004}. Starting with one for
$Z_0$, graph resolutions and principalizations by smooth-center blowups
give a K\"ahler dominating source. K\"ahlerness is preserved by these
projective modifications, and proper meromorphic images remain in
class $\mathcal C$ \cite[Lemmas~4.4 and~4.6\textup{(3)}]{FujikiDouady}.
Thus the bases also admit K\"ahler models. The preceding dimension
argument still refers to the original $Z_0$.

\subsection{Keeping the original base fixed}

\begin{corollary}[Simultaneous model control]\label{setup:simultaneous}
Suppose a smooth pair $(X_1,\Delta_1)$ has a fibration $f_1:X_1\to Y$
to a fixed smooth base, and every prime divisor contracted by $f_1$ into
codimension at least two has coefficient one in $\Delta_1$. Put
$C=B(f_1,\Delta_1)$. Let a factorization through another smooth compact
space be given, after an allowed source modification if necessary, by
\[
 X_1\xrightarrow{g_1}W_1\xrightarrow{h_1}Y.
\]
Apply \Cref{setup:flattening} to $g_1$ and make any finite number of
subsequent modifications of $W_1$, always taking dominating smooth
source models with reduced added exceptional boundaries. The resulting
diagram $X\xrightarrow{g}W\xrightarrow{h}Y$, with source boundary
$\Delta$, can be chosen to have all the following properties:
\begin{enumerate}[label=\textup{(\roman*)}]
\item Every prime contracted by $g$ into codimension at least two is
exceptional over the chosen smooth reference source for $g_1$.
\item Every prime contracted by $f=h\circ g$ into codimension at least
two in $Y$ has coefficient one in $\Delta$.
\item $B(f,\Delta)=C$. In particular, for each prime $D\subset Y$ and
each prime $E\subset X$ dominating $D$,
\begin{equation}\label{setup:retained-multiplicity}
 \ord_E(f^*D)\,m_\Delta(E)\geq m_C(D).
\end{equation}
\item The log Kodaira dimensions of the total source and of its very
general fibers over $Y$ are unchanged by these source modifications.
\end{enumerate}
The hypothesis on contracted divisors follows from
\Cref{setup:flattening}, applied to the original map to its resolved
base before fixing $Y$.
\end{corollary}

\begin{proof}
Assertion (i) is \Cref{setup:flattening}. For (ii), a source prime
is either newly exceptional, with coefficient one, or the strict
transform of an old prime with unchanged image in $Y$. The hypothesis
handles the latter case. Assertions (iii) and (iv) are
\Cref{setup:modification-compatibility};
\eqref{setup:retained-multiplicity} is the defining minimum for $C$.
Flattening the original map supplies the initial coefficient-one
condition by making its contracted primes exceptional.
\end{proof}

In applications we include the critical values and images of boundary
strata among the base bad loci, so that $C$ has SNC support. We then
fix $Y$ and use \Cref{setup:simultaneous} for changes of intermediate
bases.

\subsection{The relative Iitaka map}

\begin{lemma}[Relative Iitaka construction]
\label{application:relative-iitaka}
Let \(f:(X,\Delta)\to Y\) be a fibration between smooth compact
manifolds in Fujiki class \(\mathcal C\), with \(\Delta\) a rational
SNC boundary in \([0,1]\). Suppose that
\[
 k=\kappa(F,K_F+\Delta_F)\geq 0
\]
on very general smooth fibers.  After modifications of the source, adding
the reduced exceptional divisor to the strict transform of the boundary,
there is a factorization
\[
 X\xrightarrow{g}W\xrightarrow{h}Y
\]
with smooth compact intermediate space in class \(\mathcal C\), both
maps proper with connected fibers, and
\[
 \dim W_y=k,\qquad
 \kappa(G,K_G+\Delta_G)=0
\]
for very general fibers of \(h\) and \(g\), respectively.  The factorization
restricts to the log Iitaka fibration on very general \(f\)-fibers.
These properties persist on higher models obtained by the same rule for
the source boundary.
\end{lemma}

\begin{proof}
Consider all degrees clearing the denominators of \(\Delta\).
Their relative adjoint direct images are coherent. First shrink to
the dense analytic open where the map and boundary strata are smooth.
There the family and its adjoint line bundles are flat over the base;
generic constancy of the fiber section dimension permits cohomology
and base change \cite[\S7, nos.~4--6, S\"atze~3--5]{GrauertDirectImages},
with its corrigendum \cite{GrauertDirectImagesCorr}.
Together with the evaluation ranks this gives a dense analytic open
for each degree;
outside the resulting countable collection of exceptional sets, the
maximum image dimension is \(k\), attained in one fixed degree.

Choose a sufficiently divisible degree giving the stabilized maps,
resolve the relative meromorphic map to the projective bundle of its
direct image, and take Stein factorization onto the image. The Iitaka
fibration theorem for line bundles on compact complex spaces identifies
the resulting map on very general fibers
\cite[Chapter~II, \S5; Chapter~III, \S7]{UenoClassification}.
No finite generation is required. The image has relative dimension \(k\); its
Stein factor is in class \(\mathcal C\), being a proper image of a space
in that class.  Resolve it and the resulting source map.  Connectedness
of the fibers of \(h\) follows from that of \(f\), since their images
under \(g\) are precisely the fibers of \(h\).

The assertion that the restriction has Kodaira dimension zero concerns
the original line bundle. Its standard proof applies here: coherent
direct images on the projective Iitaka image acquire sections after an
ample twist. If the
restriction had positive Iitaka dimension, multiplying those sections
by powers of the defining system would enlarge the Iitaka image.
A nonzero defining section restricts nontrivially to a very general
Iitaka fiber, so the restriction has dimension exactly zero.

Apply this to the resolved log adjoint on a very general \(f\)-fiber.
Its divisible section ring is unchanged by
\Cref{setup:log-modification}, and adjunction identifies the restriction
with \(K_G+\Delta_G\). Further source modifications preserve the same
section ring, hence the relative Iitaka factorization and its asserted
fiber dimensions.
\end{proof}
